\documentclass[a4paper,11pt]{article}
\usepackage{amsmath,amssymb,amsthm, tikz,titlesec,hyperref,esint,braket, graphicx, mathrsfs, enumitem}
\usepackage[a4paper,margin=2cm]{geometry}
\linespread{1.3}
\newtheorem{claim}{Claim}[section]

\newcommand\name{Park Chanwoo}   % Name of the student
\newcommand\university{KAIST} % Name of the university
\newcommand\department{Physics} % Name of the department
\newcommand\studentid{20230297} % Student ID
\newcommand\s{\,\;}


\newenvironment{solution}[1]
  {\renewcommand\qedsymbol{$\square$}\begin{proof}[\textbf{Solution#1}]}
  {\end{proof}}
\newenvironment{note}
  {\renewcommand\qedsymbol{$\blacksquare$}\begin{proof}[\textnormal{\textbf{note}}]}
  {\end{proof}}

\title{KAIST\\2026 Special Topics in Physics Quantum Electronic Devices I: Interferometers and Fractional Particles\\
Homework 2\bigskip}
\author{\textbf{\Large \name} \\
% University: \university\\
Department: \department\\
Student ID: \studentid}
\date{\today}

\begin{document}
\thispagestyle{empty}
\maketitle
\tableofcontents
\titleformat{\section}[frame]{\pagebreak}{\filright
\footnotesize  \enspace \textsf{KAIST --- Interferometers and Fractional Particles 2026 Spring}\enspace}{6pt}{\Large\bfseries\filcenter}
\setcounter{secnumdepth}{1}
\setcounter{tocdepth}{2}

\newcommand{\tr}{\operatorname{tr}}
\newcommand{\sgn}{\operatorname{sgn}}
\newcommand{\supp}{\operatorname{supp}}
\renewcommand{\Re}{\operatorname{Re}}
\renewcommand{\Im}{\operatorname{Im}}
\newcommand{\boltz}{k_{\mathrm{B}}}

\section{\#1}
\subsection{(1a)}

\begin{align}
    U_{\text{BS2}}U_{\text{arm}}U_{\text{BS1}}\ket{S_1}
    &= \frac{1}{\sqrt{2}}U_{\text{BS2}}U_{\text{arm}}\left(\ket{A} + \ket{A'}\right) \\
    &= \frac{1}{\sqrt{2}}U_{\text{BS2}}\left(c_1\ket{A} + c_2\ket{A'}\right) \\
    &= \frac{1}{2}(c_1 + c_2)\ket{D_1} + \frac{1}{2}(c_1 - c_2)\ket{D_2}\\
    &\sim \cos \theta \ket{D_1} + i\sin\theta \ket{D_2}
\end{align}
Where $\theta = k\Delta = k(L_1 - L_2)$.

The probability of detecting the particle at $D_1$ is given by
\begin{equation}
    p_{D_1} = \frac{1}{4}|c_1 + c_2|^2 = |\cos \theta|^2
\end{equation}
Hence,
\begin{equation}
    \max[p_{D_1}] = 1, \quad \min[p_{D_1}] = 0
\end{equation}
and the visibility of the oscillation is
\begin{equation}
    V = \frac{\max[p_{D_1}] - \min[p_{D_1}]}{\max[p_{D_1}] + \min[p_{D_1}]} = 1
\end{equation}

\subsection{(1b)}

Inculding the effect of the environment, we have

\begin{align}
    U_{\text{BS2}}U_{\text{arm}}^{SE}U_{\text{BS1}}\ket{S_1}\ket{0_\text{env}}
    &= \frac{1}{\sqrt{2}}U_{\text{BS2}}U_{\text{arm}}^{SE}\left(\ket{A}\ket{0_\text{env}} + \ket{A'}\ket{0_\text{env}}\right) \\
    &= \frac{1}{\sqrt{2}}U_{\text{BS2}}\left(c_1\cos\varphi\ket{A}\ket{0_\text{env}} + c_1\sin\varphi\ket{A}\ket{1_\text{env}} + c_2\ket{A'}\ket{0_\text{env}}\right) \\
    &= \frac{1}{2}\ket{D_1}((c_1\cos\varphi + c_2)\ket{0_\text{env}} + c_1\sin\varphi \ket{1_\text{env}})\\ 
    &+ \frac{1}{2}\ket{D_2}((c_1\cos\varphi - c_2)\ket{0_\text{env}} + c_1\sin\varphi \ket{1_\text{env}})
\end{align}

The probability of detecting the particle at $D_1$ is changed due to the entanglement with the environment, and is given by
\begin{align}
    p_{D_1} &= \frac{1}{4}\left|c_1\cos\varphi + c_2\right|^2 + \frac{1}{4}|c_1|^2\sin^2\varphi \\
    &= \frac{1}{2}+ \frac{1}{2}\cos k\Delta\cos\varphi
\end{align}
The maximum and minimum values of $p_{D_1}$ are
\begin{equation}
    \max[p_{D_1}] = \frac{1 + |\cos\varphi|}{2}, \quad \min[p_{D_1}] = \frac{1 - |\cos\varphi|}{2}
\end{equation}
and the visibility of the oscillation is
\begin{equation}
    V = \frac{\max[p_{D_1}] - \min[p_{D_1}]}{\max[p_{D_1}] + \min[p_{D_1}]} = |\cos\varphi|
\end{equation}


\subsection{(1c)}


The expectation of the projector on $\ket{1_\text{env}}$ is
\begin{equation}
   \langle \psi |1_\text{env}\rangle\langle1_\text{env}|\psi\rangle 
    = \left\|\frac{1}{2} c_1\sin\varphi\ket{D_1} + \frac{1}{2} c_1\sin\varphi\ket{D_2} \right\|^2 = \frac{\sin^2\varphi}{2}
\end{equation}


High $\sin^2\varphi$ means that the entangement between the system and the environment is strong, and the which-path information is more accessible. Since $\sin^2\varphi + \cos^2\varphi = 1$, there is negative correlation between the which-path information and the visibility of the interference pattern.




\end{document}
